%----------------------------------------------------------------------------------------
%  Packages And Other Document Configurations
%----------------------------------------------------------------------------------------

\documentclass{resume} % Use the custom resume.cls style

% Document margins
\usepackage[left=0.75in,top=0.6in,right=0.75in,bottom=0.6in]{geometry}

% Color and hyperlink packages
\usepackage{xcolor}
\usepackage{hyperref}
\hypersetup{
  colorlinks=true,
  urlcolor=navyblue,
  linkcolor=navyblue,
  pdfborder={0 0 0}
}

% Footnote and margin adjustment packages
\usepackage{footnote}
\usepackage{changepage}

% Fontawesome package for icons
\usepackage{fontawesome}

% Tabularx package for custom tables
\usepackage{tabularx}

% Define navyblue color
\definecolor{navyblue}{RGB}{0,54,123}

%----------------------------------------------------------------------------------------
%   Customizations
%----------------------------------------------------------------------------------------

% Define italicitem, bolditem, and plainitem commands
\newcommand{\italicitem}[1]{\item{\textit{#1}}}
\newcommand{\bolditem}[1]{\item{\textbf{#1}}}
\newcommand{\plainitem}[1]{\item{#1}}

% Define user-friendly link command for hyperlinks
\newcommand{\link}[2]{{\href{#1}{#2}}}

\newcommand{\entry}[2]{#1 & #2 \tabularnewline} % Defines an entry with two arguments

%----------------------------------------------------------------------------------------
%   Define envsection command for defining a new environment section
%----------------------------------------------------------------------------------------

\newcommand{\tableEnv}[2]{%
  \begin{rSection}{#1} 
    \begin{adjustwidth}{0.0in}{0.1in} 
      \begin{tabularx}{\linewidth}{@{} >{\bfseries}l @{\hspace{6ex}} X @{}}
        #2 
      \end{tabularx}
    \end{adjustwidth}
  \end{rSection}
}

%----------------------------------------------------------------------------------------
%   Begin document
%----------------------------------------------------------------------------------------

\name{\color{navyblue} Sarang Nagpal}

\begin{document}

\printPersonalInfo{
  \personalInfo{\tag{E-mail}\info{sarangnagpal0@gmail.com}  \tag{Mobilfunknummer}\info{(+49) 175 5014497}}
  \personalInfo{\tag{Webseite}\link{https://crudybagger.github.io}{crudybagger.github.io}   \tag{Standort}\info{Aachen, Germany}}
}

%----------------------------------------------------------------------------------------
%   Summary section
%----------------------------------------------------------------------------------------

\begin{rSection}{Zusammenfassung}
Ergebnisorientierter Ingenieur und aktueller Masterstudent mit umfassender Erfahrung im Aufbau von Hochdurchsatz-Datensystemen, skalierbaren Architekturen und belastbarer KI-Infrastruktur. Nachgewiesene Fähigkeit, generative KI-Modelle, Deep-Learning-Frameworks und moderne Tools zur Lösung realer geschäftlicher Herausforderungen zu nutzen, einschließlich der Verwaltung von KI-Inferenz-Engines, die täglich mehr als eine Milliarde Transaktionen abwickeln. Hochgradig versiert im Systemdesign, unter Verwendung moderner DevOps-Praktiken, Orchestrierung von Cloud-nativen Bereitstellungen und strategischer Anwendung von Datenanalysen zur Lösung sowohl interner als auch externer Geschäftsprobleme auf Unternehmensebene.
\end{rSection}

%----------------------------------------------------------------------------------------
%   Work experience section
%----------------------------------------------------------------------------------------

\begin{rSection}{Berufserfahrung}

    \begin{rSubsection}{Visa Inc.}{April 2025 - August 2026}{Data Engineer}{Bangalore, India}
        \item Entwickelt wurde ein Token-Bucket-basierter Ratenbegrenzer in einem leistungsstarken Go-basierten verteilten System und Lastgenerator, der in Kubernetes eingesetzt wird, und ermöglicht die präzise Steuerung des Stresstestdurchsatzes bei mehr als 100.000 Transaktionen pro Sekunde, um die Systemskalierbarkeit zu verbessern.
        \item Initiierte und etablierte grundlegende Praktiken für das Chaos-Engineering in KI-Inferenzdiensten, die für die Risikomodellierung von Finanztransaktionen verwendet werden, wobei der Schwerpunkt auf der proaktiven Injektion von Fehlern in die Kubernete und Vanille-Cloud-Infrastruktur liegt, um die Systemresilienz zu validieren.
        \item Entwickelte Chaos-Testfälle für mehrere Komponenten der Architektur, die Latenz, Ressourcenerschöpfung und veraltete Kontextszenarien simulieren, um die Robustheit der Risk Modeling AI-Inferenz-Engine zu testen.
        \item Entwickelte ein Tool-agnostisches Jenkins-basiertes Cron-System, um die Ausführung von Chaos-Tests zu automatisieren und detaillierte Berichte über die Systemleistung zu erstellen, die eine kontinuierliche Verbesserung der KI-Inferenzzuverlässigkeit ermöglichen.
        \item Zusammenarbeit funktionsübergreifend mit Datenwissenschaftlern und Architekten, um Telemetriemetrie- und Observability-Dashboards in Prometheus und Grafana zu entwickeln und den Explosionsradius von Chaostests bei KI-Inferenzen genau zu messen.
        \item Optimierte und fein abgestimmte Kubernetes verteilten Golang- und Kafka-Datenpipelines, um den Durchsatz zu verbessern, um sicherzustellen, dass das System strenge SLAs für die Verarbeitung von Finanztransaktionen in Echtzeit erfüllt.                  
        \item Verbesserte Plattformverfügbarkeit hin zu einem 99,99999\%ile-Ziel, indem sie zur Verteilung der Infrastruktur in unabhängige betriebliche Stapel beiträgt, die über mehrere Availability Zones verteilt sind, um eine hohe Verfügbarkeit und Notfallwiederherstellung für kritische Finanzdienstleistungen zu gewährleisten.
    \end{rSubsection}
    
    \begin{rSubsection}{Amazon}{January 2024 - February 2025}{Software Development Engineer}{Bangalore, India}
        \item Verbesserte die Onboarding-Erfolgsrate der Unified Payments Interface (UPI)-Zahlungsmethode um 5 \% für eine aktive Basis von über 20 Millionen Kunden, indem der Onboarding-Prozess verkürzt und die Anzahl der für Neuregistrierungen erforderlichen Klicks reduziert wurde.
        \item Entwickelte ein System, um bestehende UPI-Kunden über mehrere Zahlungsdienstleister hinweg sicher zu registrieren, ohne Onboarding-Dienste zu skalieren, wodurch Probleme mit einem Ausfall durch einen Punkt des Ausfalls vermieden wurden.
        \item Entwickelte und entwickelte robuste RESTful-APIs, die Java neben einer Reihe von serverlosen AWS-Komponenten, einschließlich AWS Lambda, DynamoDB, SQS und SNS, verwenden.
        \item Modernisierte ältere mobile Web-Interfaces von JavaServer Pages (JSP) bis ReactJS mit Redux, was zu einer 8\% Verbesserung der Onboarding-Metriken und einer reduzierten Latenz führte.
    \end{rSubsection}

    \begin{rSubsection}{Amazon}{May 2022 - July 2022}{Software Development Engineer Intern}{Bangalore, India}
        \item Konzeptualisiert und einen eigenständigen Backend-Service mit Spring Boot (Java) erstellt, um einen Core-UPI-Service zu testen.
        \item Skalierte EC2-Instance-Infrastruktur dynamisch, um benutzerdefinierte Stresstest-Frameworks in einer simulierten Hochsicherheitsumgebung erfolgreich durchzuführen.
        \item Erzielte eine 100 \%-Punktzahl in der Resilienzmetrik durch Messung und Dokumentation von Leistungsschwellen vor Abschluss des Praktikums.
    \end{rSubsection}
    
\end{rSection}

%----------------------------------------------------------------------------------------
%   Education section
%----------------------------------------------------------------------------------------

\begin{rSection}{Bildung}

    \begin{rSubsectionNoBullet}{\bf Master of Science}{RWTH Aachen University}{Computer Engineering}{May 2026 - Present}
        \italicitem{}
    \end{rSubsectionNoBullet}
    
    \begin{rSubsectionNoBullet}{\bf Bachelor of Technology }{National Institute of Technology, Warangal}{Electronics and Communication Engineering}{May 2019 - June 2023}
        \italicitem{Entwickelt eine starke mathematische Grundlage in Steuerungssystemen, fortschrittlicher Physik, algorithmischem Design auf niedrigem Niveau und Physik und Design von Halbleiterbauelementen.}
    \end{rSubsectionNoBullet}

\end{rSection}

%----------------------------------------------------------------------------------------
%   Skills section
%----------------------------------------------------------------------------------------

\tableEnv{Fähigkeiten}{
    \entry{Programmierung}{Go (Golang), Python, C++, Java, JavaScript, TypeScript, SQL}
    
    \entry{KI \& Datenanalyse}{Generative KI, Große Sprachmodelle (LLMs), TensorFlow, Keras, Pandas, Prompt Engineering, Datenaggregation}
    
    \entry{Cloud \& DevOps}{AWS (Lambda, DynamoDB, SQS, SNS), Kubernetes, Docker, ChaosMesh, Kafka, RabbitMQ, Jenkins}
    
    \entry{System Design}{Microservices, Serverless Architecture, Event-Driven Messaging, RESTful APIs, Distributed Systems, High-Throughput Data Pipelines}

    \entry{Languages}{English (C2 - IELTS 8.0), German (A2 - Goethe Zertifikat)}
}

%----------------------------------------------------------------------------------------
%   Certifications section
%----------------------------------------------------------------------------------------

\begin{rSection}{Certifications}

    \begin{rSubsectionNoBullet}{\bf Certified Kubernetes Application Developer (CKAD)}{The Linux Foundation}{Cloud-Native Architecture Certification}{August 2026}
        \plainitem{Praktische Kompetenz bei der Gestaltung, Erstellung, Konfiguration und Bereitstellung komplexer Cloud-nativer Anwendungen direkt für verteilte Kubernetes-Umgebungen.}
    \end{rSubsectionNoBullet}
    
    \begin{rSubsectionNoBullet}{\bf Machine Learning Certification}{Stanford University (Coursera)}{Algorithmische Stiftung Zertifizierung}{December 2022}
        \plainitem{Mathematische Grundlage des überwachten und unüberwachten Lernens.}
    \end{rSubsectionNoBullet}
    
    \begin{rSubsectionNoBullet}{\bf Deep Learning Specialization}{Andrew Ng (Coursera)}{Angewandte Neurale Netzwerke Zertifizierung}{September 2020}
        \plainitem{Convolutional Neural Networks (CNNs); Sequence Models; hyperparameter tuning; mathematical regularization; optimizing deep networks.}
    \end{rSubsectionNoBullet}

\end{rSection}

\end{document}
